% Copyright 2024 Richard J. Zak
% richard.j.zak@gmail.com

\documentclass[letter,10pt]{article}
\usepackage[breaklinks]{hyperref}
\hypersetup{
    bookmarks=true,         % show bookmarks bar?
    unicode=false,          % non-Latin characters in Acrobat’s bookmarks
    pdftoolbar=true,        % show Acrobat’s toolbar?
    pdfmenubar=true,        % show Acrobat’s menu?
    pdffitwindow=false,     % window fit to page when opened
    pdfstartview={XYZ null null 1.00},    % disable zoom
    pdftitle={Practice Quiz 3},    % title
    pdfauthor={Richard Zak},     % author
    pdfsubject={UMBC CMSC104 Problem Solving and Computer Programming},   % subject of the document
    pdfkeywords={Computer Science, Programming, Problem Solving, CSEE}, % list of keywords
    pdfnewwindow=true,      % links in new PDF window
    colorlinks=false,       % false: boxed links; true: colored links
    linkcolor=red,          % color of internal links (change box color with linkbordercolor)
    citecolor=green,        % color of links to bibliography
    filecolor=magenta,      % color of file links
    urlcolor=cyan           % color of external links
}
\usepackage{graphicx}
\usepackage{fancyhdr}
\usepackage{multicol}
\pagestyle{fancy}
\usepackage[letterpaper, margin=1in]{geometry}
\geometry{letterpaper}
\usepackage{listings} % Syntax highlighing
\usepackage{amsmath}
\usepackage{parskip} % Disable initial indent
\usepackage{color,soul} % Highligher
\usepackage[normalem]{ulem} % Strikethrough with \sout{}

\definecolor{codegreen}{rgb}{0,0.6,0}
\definecolor{codegray}{rgb}{0.5,0.5,0.5}
\definecolor{codepurple}{rgb}{0.58,0,0.82}
\definecolor{backcolour}{rgb}{0.97,0.97,0.97}

\lstdefinestyle{mystyle}{
    backgroundcolor=\color{backcolour},
    commentstyle=\color{codegreen},
    keywordstyle=\color{magenta},
    numberstyle=\tiny\color{codegray},
    stringstyle=\color{codepurple},
    basicstyle=\ttfamily\small,
    breakatwhitespace=false,
    breaklines=true,
    captionpos=b,
    keepspaces=true,
    numbers=left,
    numbersep=5pt,
    showspaces=false,
    showstringspaces=false,
    showtabs=false,
    tabsize=2
}

\lstset{style=mystyle}

\usepackage[utf8]{inputenc}
\fancyhf{}
\renewcommand{\headrulewidth}{0pt} % Remove default underline from header package
\rhead{CMSC 104 Section 01
: Practice Quiz 3}
%\rhead{}
\lhead{\begin{picture}(0,0) \put(0,-10){\includegraphics[width=1.1cm]{../Images/UMBC-vertical}} \end{picture}}
\cfoot{\thepage}
\rfoot{Fall 2026
}
\lfoot{CMSC 104 Section 01
}
\AtEndDocument{\vfill \footnotesize{Last modified: 20 October 2026}}
\AtEndDocument{\rfoot{Fall 2026
}}
\renewcommand\thesubsection{\arabic{subsection}} % Show only subsection numbers, not section.subsection

\begin{document}

\huge
\textbf{Practice Quiz 3}
\normalsize

\paragraph{}Quiz 3 covers chapters 1 through 9. Please answer the following questions. Partial credit may be given for incomplete free response questions.

\begin{enumerate}
	\item Which of the following item types are mutable?
    \begin{itemize}
        \item List
        \item Dictionary
        \item Tuple
        \item Class
        \item Set
        \item String
    \end{itemize}

    \item A sentinel controlled loop is a type of event loop. But what about it makes it different? What are any limitations?

    \item What are examples for when you'd use a \texttt{for} loop vs. a \texttt{while} loop?

    \item What are the types of loops discussed in class?
    \begin{enumerate}
        \item Definitive loops (counter-controlled loops)
		\item Indefinite loops (event-controlled loops)
		\item Gaming loop (optimized for compute-intensive tasks, such as video games)
		\item Even-odd loop (skips odd iterations)
		\item Operating System loops (used only by operating systems to keep the computer powered on)
    \end{enumerate}

    \item Given: \verb|a = True, b = False|, evaluate the following boolean instructions and indicate if they're true or false.
    \begin{enumerate}
        \item \verb|a and not b == (a or not(a))|
        \item \verb|not (a or b) == not(a) and not(b)|
        \item \verb|b == 0|
        \item \verb|a and not(b) == a and b|
    \end{enumerate}

    \item Accessing a single character in a string is called what?
    \begin{enumerate}
        \item Slicing
        \item Indexing
        \item Concatenating
        \item Assigning
    \end{enumerate}

    \item Merging two or more strings into a larger string is called what?
    \begin{enumerate}
        \item Slicing
        \item Indexing
        \item Concatenating
        \item Assigning
    \end{enumerate}

    \item Getting a sub-string from a larger string ("science" from "Computer science is fun") is called what?
    \begin{enumerate}
        \item Slicing
        \item Indexing
        \item Concatenating
        \item Assigning
    \end{enumerate}

    \item A list must contain a minimum of one item.

    \item A list or a tuple must contain elements of the same type.

    \item If you write an if-else statement with a nested if-else statement and find that you have blocks with identical code, what might be done differently?
\end{enumerate}

\end{document}